\chapter{Asymptotic expansions and summation estimates}
\label{module:M03.17}
\SourceRef{Zorich I, Appendix D; Zorich II, chapter 19; selected real-variable
parts of Olver, chapters 1, 3 and 8}

\KEntry{K03.17.asymptotic_scale}{Scales and finite expansions}
\begin{definition}
Along a specified limit, an asymptotic scale is a family $\phi_0,\phi_1,\ldots$
of eventually nonzero scalar functions such that $\phi_{j+1}=o(\phi_j)$.
An expansion through index $m$ means the finite assertion
\[
 f=\sum_{j=0}^{m}a_j\phi_j+o(\phi_m).
\]
An all-orders expansion means this assertion for every finite $m$, with
consistent coefficients. It does not assert convergence of an infinite sum.
\end{definition}
For $j<k$, $\phi_k/\phi_j$ is a finite product of adjacent ratios tending to
zero and hence tends to zero. The domain and direction of the limit are part
of the definition; a two-sided real power scale often needs $|x|$ in estimates.

\KEntry{K03.17.coefficient_unique}{Uniqueness does not determine the function}
\begin{proposition}
The coefficients of an expansion in a fixed asymptotic scale are unique.
Functions can nevertheless have the same all-orders expansion.
\end{proposition}
\begin{proof}
Subtract two expansions of the same function and let $j$ be the least index
with unequal coefficients. Divide by $\phi_j$. All terms of larger index
tend to zero, as does the remainder divided by $\phi_j$; the nonzero constant
difference at index $j$ cannot tend to zero. This is a contradiction.
Uniqueness concerns coefficients of a given function, not uniqueness of that
function; \KRef{K03.17.flat_remainder} gives the latter distinction.
\end{proof}

\KEntry{K03.17.operations}{Arithmetic with controlled remainders}
\begin{proposition}
Finite asymptotic expansions can be added and multiplied by retaining the
resulting remainders. Division is valid after an eventually nonzero leading
term has been factored out. Closure of an arbitrary scale under multiplication
or composition is not automatic.
\end{proposition}
\begin{proof}
Write $f=P+r$, $g=Q+s$, where $P,Q$ are the stated finite sums. The exact
identities
\[
 f+g=P+Q+r+s,\qquad fg=PQ+Ps+Qr+rs
\]
reduce every truncation to the $O/o$ multiplication rules. For example, if
$f=\sum_{j=0}^m a_jt^j+o(t^m)$ and
$g=\sum_{j=0}^m b_jt^j+o(t^m)$ as $t\to0$, then $P,Q=O(1)$;
the coefficient at $t^k$ is $\sum_{j=0}^ka_jb_{k-j}$ for $k\leq m$, and
the discarded finite polynomial terms of degree $>m$ and the three remainder
terms are $o(t^m)$. For a nonzero leading coefficient, divide first by the
leading scale factor and write the remaining denominator as $a(1+u)$,
$a\ne0$, $u\to0$. The finite identity
\[
 \frac1{1+u}=\sum_{j=0}^{m}(-u)^j+
 \frac{(-u)^{m+1}}{1+u}
\]
both proves eventual invertibility and controls the quotient remainder.
Composition requires the inner variable to approach the expansion point and
its range to remain in the outer domain. For instance $u(t)=O(t)$ and an outer
remainder $o(u^m)$ imply a remainder $o(t^m)$ by the defining epsilon bound,
including points with $u=0$ after the remainder is assigned its correct value.
For a nonpower scale one must keep its product scales or explicitly justify
their order before dropping terms.
\end{proof}

\KEntry{K03.17.flat_remainder}{Beyond all powers}
\begin{example}
The function $e^{-1/x^2}$ for $x\ne0$, extended by zero at $0$, is
$o(|x|^m)$ for every nonnegative integer $m$, but is positive off $0$.
Thus its all-orders power expansion is the zero expansion without the function
being zero on a neighborhood.
\end{example}
\begin{proof}
Put $u=1/x^2$. For an integer $r>m/2$, the exponential series gives
$e^u\geq u^r/r!$ for $u>0$, so
$u^{m/2}e^{-u}\leq r!u^{m/2-r}\to0$. This proves the claimed estimates.
The full smooth-flat construction, including all derivatives at zero, is
\KRef{K03.06.smooth_flat}; no analyticity inference is made here.
\end{proof}

\KEntry{K03.17.divergent_asymptotic}{A divergent series that is asymptotically correct}
\begin{example}
For $x>0$ let $F(x)=\int_0^\infty e^{-t}/(1+xt)\dd t$. For every integer
$N\geq1$,
\[
 F(x)=\sum_{k=0}^{N-1}(-1)^kk!x^k+R_N(x),\qquad
 |R_N(x)|\leq N!x^N.
\]
The associated infinite power series has radius zero.
\end{example}
\begin{proof}
The integral is absolutely convergent by comparison with $e^{-t}$. Insert
the exact finite geometric identity into its integrand. Each moment is
$\Gamma(k+1)=k!$, and the remainder has absolute value at most
$x^N\int_0^\infty t^Ne^{-t}\dd t=N!x^N$, since $1+xt\geq1$.
For fixed $x\ne0$, consecutive absolute power-series terms have ratio
$(k+1)|x|\to\infty$, so they do not even tend to zero. The finite estimates
remain valid and, using one extra term, give an all-orders asymptotic expansion.
\end{proof}

\KEntry{K03.17.bernoulli}{Bernoulli polynomials and their convention}
\begin{definition}
Set $B_0(x)=1$. Recursively, $B_n$ is the unique polynomial satisfying
$B_n'=nB_{n-1}$ and $\int_0^1B_n(x)\dd x=0$ for $n\geq1$.
Write $B_n=B_n(0)$ for the Bernoulli number. In particular,
$B_1(x)=x-1/2$, so $B_1=-1/2$.
\end{definition}
\begin{proof}[Construction and identities]
Integrate the already constructed polynomial $nB_{n-1}$ and choose the unique
constant giving zero average. This proves existence and uniqueness. It gives
$B_2(x)=x^2-x+1/6$ and $B_4(x)=x^4-2x^3+x^2-1/30$ by coefficient comparison.
For $n\geq2$,
$B_n(1)-B_n(0)=n\int_0^1B_{n-1}=0$.
The polynomials $(-1)^nB_n(1-x)$ satisfy the same derivative and average
recursion, so uniqueness gives $B_n(1-x)=(-1)^nB_n(x)$.
Consequently $B_n(0)=0$ for odd $n\geq3$.

Induction on $n$, by differentiating and then evaluating at $x=0$, gives
$B_n(x+1)-B_n(x)=nx^{n-1}$; the case $n=1$ is direct and the endpoint
identity fixes the induction constant for $n\geq2$. Telescoping therefore
proves the finite power-sum formula
\[
 \sum_{k=0}^{N-1}k^p=
 \frac{B_{p+1}(N)-B_{p+1}(0)}{p+1}\quad(p\geq0).
\]
For $p=0$ interpret the summands as $1$, including the $k=0$ term.
\end{proof}

\KEntry{K03.17.euler_maclaurin}{Euler--Maclaurin with an actual remainder}
\begin{theorem}
Let integers $a<b$, $m\geq1$, and $f\in C^{2m}([a,b],\R\text{ or }\C)$.
Then
\begin{align*}
 \frac{f(a)}2+\sum_{a<k<b}f(k)+\frac{f(b)}2
 &=\int_a^b f(x)\dd x+
 \sum_{r=1}^m\frac{B_{2r}}{(2r)!}
 \bigl(f^{(2r-1)}(b)-f^{(2r-1)}(a)\bigr)\\
 &\quad-\frac1{(2m)!}\int_a^b
 B_{2m}(\{x\})f^{(2m)}(x)\dd x,
\end{align*}
where $\{x\}=x-\lfloor x\rfloor$.
\end{theorem}
\begin{proof}
Integrating $(x-k-1/2)f'(x)$ by parts on $[k,k+1]$ gives
\[
 \frac{f(k)+f(k+1)}2-\int_k^{k+1}f
 =\int_k^{k+1}B_1(x-k)f'(x)\dd x.
\]
Sum over the unit intervals. Put
$T_j=\int_a^b B_j(\{x\})f^{(j)}(x)\dd x/j!$.
For $1\leq j<2m$, integrate by parts on each unit interval and sum.
The equality $B_{j+1}(1)=B_{j+1}(0)$ cancels internal boundary terms, yielding
\[
 T_j=\frac{B_{j+1}}{(j+1)!}
 \bigl(f^{(j)}(b)-f^{(j)}(a)\bigr)-T_{j+1}.
\]
Odd Bernoulli numbers beyond $B_1$ vanish. Starting from $T_1$ and applying
this identity successively leaves exactly the even-index corrections and
$-T_{2m}$ in the theorem. The periodic polynomials have at most finitely
many relevant junctions on $[a,b]$; their endpoint values at those junctions
do not change their proper Riemann integrals. The stated regularity permits
every integration by parts used above.
\end{proof}

\KEntry{K03.17.em_bound}{A useful finite error bound}
\begin{corollary}
With $M_{2m}=\max_{0\leq t\leq1}|B_{2m}(t)|$, the absolute remainder is
at most
\[
 \frac{M_{2m}}{(2m)!}\int_a^b|f^{(2m)}(x)|\dd x.
\]
\end{corollary}
\begin{proof}
Take the absolute value inside the remainder integral and use the uniform
polynomial bound $M_{2m}$. This finite bound supplies no convergence
claim for an infinite Euler--Maclaurin series.
\end{proof}

\KEntry{K03.17.harmonic}{Euler's constant and harmonic corrections}
\begin{proposition}
There is a real constant $\gamma$ such that
\[
 H_n=\log n+\gamma+\frac1{2n}-\frac1{12n^2}
       +\frac1{120n^4}+O(n^{-6}),\qquad H_n=\sum_{k=1}^n\frac1k.
\]
More generally each fixed number of Bernoulli corrections has a controlled
inverse-power remainder.
\end{proposition}
\begin{proof}
Integral comparison gives $H_n\geq\log(n+1)>\log n$ and
$\log(1+1/n)>1/(n+1)$; hence $H_n-\log n$ is positive decreasing and has
a limit $\gamma$. Apply Euler--Maclaurin to $f(x)=1/x$ on $[n,M]$.
Its $(j)$th derivative is $(-1)^jj!x^{-j-1}$. The endpoint-half sum is
$H_M-H_n+1/(2n)-1/(2M)$. Sending $M\to\infty$ gives
\[
 H_n=\log n+\gamma+\frac1{2n}
 -\sum_{r=1}^m\frac{B_{2r}}{2r\,n^{2r}}+E_m(n),
 \qquad |E_m(n)|\leq\frac{M_{2m}}{2m\,n^{2m}}.
\]
The bound follows by integrating $x^{-2m-1}$ and justifies the improper
limit of the remainder. Taking the formula with $m=3$, and absorbing its
last correction and remainder into $O(n^{-6})$, proves the displayed
formula, since $B_2=1/6$ and $B_4=-1/30$. In general one more correction
than those displayed gives the desired order of error.
\end{proof}

\KEntry{K03.17.stirling}{Stirling's formula, including the constant}
\begin{theorem}
As $n\to\infty$ through positive integers,
\[
 n!=\sqrt{2\pi n}(n/e)^n\exp(O(1/n))
     \sim\sqrt{2\pi n}(n/e)^n.
\]
\end{theorem}
\begin{proof}
Put $s_n=\log(n!)-(n+1/2)\log n+n$. A direct subtraction gives
\[
 s_{n+1}-s_n=1-(n+1/2)\log(1+1/n)
            =-\frac1{12n^2}+O(n^{-3}).
\]
The last equality follows from the finite Taylor expansion of $\log(1+u)$
with a bounded fourth derivative on $[0,1/2]$; the finitely many smaller
indices do not matter. Thus $\sum|s_{n+1}-s_n|$ converges and its tail is
$O(1/n)$ by the $p$-series integral bound. Write $s_n=c+O(1/n)$ and
$A=e^c>0$. Exponentiation gives
$n!=A n^{n+1/2}e^{-n}\exp(O(1/n))$.

To determine $A$, form $c_n=\binom{2n}{n}/4^n$. The preceding formula yields
$\sqrt n\,c_n\to\sqrt2/A$. Independently, the Wallis calculation
\KRef{K03.18.wallis} gives $\sqrt n\,c_n\to1/\sqrt\pi$.
Uniqueness of a finite limit gives $A=\sqrt{2\pi}$. This identifies the
normalization rather than hiding it in an unspecified multiplicative constant.
The logarithmic $O(1/n)$ error also implies a relative $1+O(1/n)$ error,
by the mean-value bound for $e^u$ near zero.
\end{proof}

\KEntry{K03.17.laplace_local}{An interior Laplace minimum}
\begin{theorem}
Let $a<b$, let $\phi\in C^2([a,b],\R)$ have a unique global minimum at
$x_0\in(a,b)$, and suppose $d=\phi''(x_0)>0$. For continuous scalar $g$,
\[
 \sqrt\lambda\,e^{\lambda\phi(x_0)}
 \int_a^b e^{-\lambda\phi(x)}g(x)\dd x
 \longrightarrow g(x_0)\sqrt{\frac{2\pi}{d}}
 \quad(\lambda\to+\infty).
\]
The assertion includes $g(x_0)=0$; it is a normalized limit rather than an
invalid asymptotic equivalence to zero.
\end{theorem}
\begin{proof}
Fermat gives $\phi'(x_0)=0$. Taylor--Peano gives
$q(x_0+h):=\phi(x_0+h)-\phi(x_0)=dh^2/2+o(h^2)$.
Choose $r>0$ with $[x_0-r,x_0+r]\subset(a,b)$ and
$q(x_0+h)\geq c h^2$ there for some $c>0$, for instance $c=d/4$.
On the complementary compact set uniqueness of the minimum gives
$q\geq\eta>0$. If $M=\max|g|$, its normalized contribution is at most
$\sqrt\lambda M(b-a)e^{-\lambda\eta}\to0$.

In the local integral substitute $y=\sqrt\lambda(x-x_0)$. The normalized
integrand becomes
$g(x_0+y/\sqrt\lambda)e^{-\lambda q(x_0+y/\sqrt\lambda)}$ on
$|y|\leq r\sqrt\lambda$, bounded in absolute value by $Me^{-cy^2}$.
For fixed $R$, the defining $o(h^2)$ bound is uniform for
$|h|\leq R/\sqrt\lambda$ and proves uniform convergence on $[-R,R]$ to
$g(x_0)e^{-dy^2/2}$. The tails outside this fixed interval are bounded by
$M\int_{|y|>R}e^{-cy^2}\dd y$, uniformly in large $\lambda$.
They tend to zero as $R\to\infty$, by the Gaussian integral and a linear
substitution. First choose $R$ for the tails, then $\lambda$ for the proper
integrals. This proves convergence to
$g(x_0)\int_\R e^{-dy^2/2}\dd y=g(x_0)\sqrt{2\pi/d}$.
The argument uses explicit proper convergence and tails, not an unstated
dominated-convergence theorem.
\end{proof}
\begin{corollary}[A quantitative source variant]
If additionally $\phi$ is $C^3$ near $x_0$ and $g$ is locally Lipschitz there,
the normalized expression equals
$g(x_0)\sqrt{2\pi/d}+O(\lambda^{-1/2})$.
\end{corollary}
\begin{proof}
In the preceding variables,
$|\lambda q(x_0+y/\sqrt\lambda)-dy^2/2|
 \leq C|y|^3/\sqrt\lambda$ and
$|g(x_0+y/\sqrt\lambda)-g(x_0)|\leq L|y|/\sqrt\lambda$.
Both exponents are at least $c'y^2$ after reducing $r$.
The elementary inequality $|e^{-u}-e^{-v}|\leq|u-v|e^{-\min(u,v)}$
therefore bounds the local difference by
$C'\lambda^{-1/2}(|y|+|y|^3)e^{-c'y^2}$.
Its integral is finite, by substitution $t=c'y^2$ and Gamma moments.
The omitted Gaussian and exterior tails are exponentially small in
$\lambda$, proving the error estimate. These stronger hypotheses are not
attached to the weaker normalized-limit theorem.
\end{proof}

\KEntry{K03.17.watson}{A real Watson lemma with tail conditions}
\begin{theorem}
Let $N\geq1$, $\alpha>0$. Suppose $q$ is locally Riemann integrable on
$(0,\infty)$ and, for $0<t\leq\delta$,
\[
 q(t)=\sum_{k=0}^{N-1}a_kt^{\alpha+k-1}
             +O(t^{\alpha+N-1}).
\]
Suppose also $|q(t)|\leq C e^{ct}$ for all sufficiently large $t$.
Then, as $\lambda\to+\infty$,
\[
 \int_0^\infty e^{-\lambda t}q(t)\dd t
 =\sum_{k=0}^{N-1}\frac{a_k\Gamma(\alpha+k)}{\lambda^{\alpha+k}}
       +O(\lambda^{-\alpha-N}).
\]
\end{theorem}
\begin{proof}
The local expansion bounds $q$ by a constant times $t^{\alpha-1}$ near zero,
so the integral there converges absolutely. The growth bound controls the
infinite tail for $\lambda>c$. On $[\delta,A]$ local integrability gives a
finite absolute integral, bounded after damping by
$e^{-\lambda\delta}\int_\delta^A|q|$. On $[A,\infty)$ the bound is
$Ce^{-(\lambda-c)A}/(\lambda-c)$. Thus the whole tail beyond $\delta$ is
exponentially small for large $\lambda$; so are the tails of each of the
finitely many power monomials beyond $\delta$.

The local remainder contributes at most
\[
 C'\int_0^\delta e^{-\lambda t}t^{\alpha+N-1}\dd t
 \leq C'\Gamma(\alpha+N)\lambda^{-\alpha-N}.
\]
For each retained monomial, substitute $u=\lambda t$ in its full improper
integral to obtain $\Gamma(\alpha+k)\lambda^{-\alpha-k}$.
Adding the finite contributions and the controlled tails proves the formula.
The global growth condition is what prevents a formally correct local
expansion from being misused for an uncontrolled infinite integral.
\end{proof}

\KEntry{K03.17.oscillatory_bound}{An elementary oscillatory estimate}
\begin{proposition}
For $h\in C^1([a,b],\C)$ and nonzero real $\lambda$,
\[
 \left|\int_a^b h(x)e^{i\lambda x}\dd x\right|
 \leq\frac{|h(a)|+|h(b)|+\int_a^b|h'(x)|\dd x}{|\lambda|}.
\]
\end{proposition}
\begin{proof}
Real-parameter differentiation of the established complex exponential gives
$(e^{i\lambda x})'=i\lambda e^{i\lambda x}$. Integration by parts yields
the boundary term $[h(x)e^{i\lambda x}/(i\lambda)]_a^b$ and the remaining
integral $-(i\lambda)^{-1}\int h'e^{i\lambda x}$. Take absolute values and
use $|e^{i\lambda x}|=1$.
\end{proof}

\KEntry{K03.17.scope_boundary}{What this asymptotic supplement does not assert}
\begin{remark}[Routing record]
General stationary phase, steepest descent, Stokes transitions, WKB and uniform
Airy/Bessel approximations require additional D12/D13/D24 foundations and
estimates. This record routes those subjects and contains no claimed proof of
their general theories. All statements proved in this chapter are finite
real-variable estimates with explicit remainders or tail control.
\end{remark}
